% sec:dg-metric — The Metric Tensor and Index Operations
\section{The Metric Tensor and Index Operations}
\label{sec:dg-metric}

To trace a photon near a black hole, the ray tracer must answer geometric
questions at every integration step: how far apart are two nearby rays?  How
much proper time elapses between events along a timelike worldline?  Is a
given displacement timelike, null, or spacelike?  The tangent-space machinery
from \cref{sec:dg-manifolds} cannot answer any of these --- the spaces
$T_p\mathcal{M}$ and $T^*_p\mathcal{M}$ carry no notion of length, angle,
or distance.  What is missing is an inner product on each tangent space: the
metric tensor.
\coderef{Math.Tensor.Symmetric}{Sym4}

\paragraph{The metric tensor.}

A \emph{metric tensor} $\metric$ at a point $p$ is a symmetric,
non-degenerate bilinear form on $T_p\mathcal{M}$.  In local coordinates
$x^\mu$, where $V = V^\mu\,\pd{\mu}$ and $W = W^\mu\,\pd{\mu}$, its action
takes the component form
%
\begin{equation}
  g(V, W) = \metric\, V^\mu\, W^\nu \,.
  \label{eq:metric-bilinear}
\end{equation}
%
Symmetry means $g_{\mu\nu} = g_{\nu\mu}$, so in four dimensions the metric has
$4 \times 5 / 2 = 10$ independent components rather than 16.  Non-degeneracy
means the determinant $g = \det(\metric)$ is non-zero everywhere --- no
tangent vector is ``invisible'' to the metric.  In short, the metric is the
curved-spacetime generalisation of the dot product: it assigns a ``size'' to
displacements and an ``angle'' between directions.

\begin{definition}[Metric tensor]
\label{def:metric-tensor}
A \emph{Lorentzian metric} on a smooth manifold $\mathcal{M}$ is a smooth
assignment $p \mapsto g_p$ of a symmetric, non-degenerate bilinear form of
signature $(-,+,+,+)$ to each tangent space $T_p\mathcal{M}$.  The pair
$(\mathcal{M}, g)$ is a \emph{Lorentzian manifold} --- the mathematical arena
for general relativity.
\end{definition}

\implnote{The codebase represents the metric as \texttt{Sym4 'Covariant},
storing exactly 10 independent components.  The constructor
\texttt{sym4Diag (-1) 1 1 1 :: Sym4 'Covariant} produces the Minkowski
metric in Cartesian coordinates.}

\paragraph{The line element.}

The metric's most direct physical expression is the \emph{line element},
which assigns a squared interval to an infinitesimal coordinate displacement
$\d x^\mu$:
%
\begin{equation}
  \d s^2 = \metric\,\d x^\mu\,\d x^\nu \,.
  \label{eq:line-element}
\end{equation}
%
This is a compact notation for a double sum over $\mu$ and $\nu$ from 0 to 3,
with the Einstein convention summing over both repeated index pairs.  For the
Minkowski metric $\eta_{\mu\nu} = \operatorname{diag}(-1,+1,+1,+1)$ from
\cref{sec:sr-minkowski}, the line element reduces to
$\d s^2 = -\d t^2 + \d x^2 + \d y^2 + \d z^2$ --- the interval that
underpinned all of \cref{ch:special-relativity}.  A general curved-spacetime
metric replaces these constant coefficients with functions of the coordinates.
\physnote{For a timelike displacement ($\d s^2 < 0$), the proper time is
$\d\tau = \sqrt{-\d s^2}$.  For a spacelike displacement ($\d s^2 > 0$),
$\d s$ is the proper distance.  Null displacements have $\d s^2 = 0$ exactly.}

\paragraph{Signature and causal character.}

The single minus sign in the signature $(-,+,+,+)$ is what makes spacetime
\emph{spacetime} rather than a four-dimensional Riemannian manifold.  At each
point $p$, the metric divides tangent vectors into three classes by the sign of
$g(V, V) = \metric\, V^\mu\, V^\nu$:
%
\begin{equation}
  g(V, V) \;\begin{cases}
    < 0 & \text{timelike} \\
    = 0 & \text{null (lightlike)} \\
    > 0 & \text{spacelike}
  \end{cases}
  \label{eq:causal-character}
\end{equation}
%
This classification is coordinate-independent --- it depends on the metric and
the vector, not on the chart.  The causal structure introduced in
\cref{sec:sr-causal} for flat spacetime now extends to every point on the
manifold: each tangent space carries its own light cone, and the metric
determines its shape.
\xrefnote{The light cone at each point was introduced in
\cref{sec:sr-causal}; in curved spacetime, these cones tilt and deform from
point to point.}

\paragraph{Index lowering and raising.}

The metric provides the canonical map between vectors and covectors that
\cref{sec:dg-manifolds} noted was missing from the bare tangent-space
framework.  Given a tangent vector $V^\mu$, we define its \emph{associated
one-form} by
%
\begin{equation}
  V_\mu = \metric\, V^\nu \,.
  \label{eq:index-lower}
\end{equation}
%
This operation is called \emph{lowering an index}, and it maps
$T_p\mathcal{M} \to T^*_p\mathcal{M}$.  The inverse operation,
\emph{raising an index}, requires the inverse metric $\inversemetric$
(defined below):
%
\begin{equation}
  V^\mu = \inversemetric\, V_\nu \,.
  \label{eq:index-raise}
\end{equation}
%
Together, these two operations constitute the \emph{musical isomorphism} ---
lowering is denoted $\flat$ and raising $\sharp$ in differential-geometric
notation.  Non-degeneracy of the metric guarantees that the map is an
isomorphism: every vector has a unique associated covector, and vice versa.
\coderef{Math.Tensor.Operations}{lower4}

\warnnote{Lowering and raising change the object's type: $V^\mu$ lives in
$T_p\mathcal{M}$, while $V_\mu$ lives in $T^*_p\mathcal{M}$.  Even for the
Minkowski metric, the components differ: if $V^\mu = (E, p^x, p^y, p^z)$
then $V_\mu = (-E, p^x, p^y, p^z)$.}

\paragraph{The inverse metric.}

The \emph{inverse metric} $\inversemetric$ is defined by the condition
%
\begin{equation}
  g^{\mu\alpha}\, g_{\alpha\nu} = \delta^\mu{}_\nu \,,
  \label{eq:inverse-metric}
\end{equation}
%
where $\delta^\mu{}_\nu$ is the Kronecker delta.  Since $\metric$ is
non-degenerate, the matrix $(g_{\mu\nu})$ is invertible, and $\inversemetric$
is its matrix inverse.  The inverse metric is itself a symmetric $(2,0)$
tensor --- it takes two one-forms and produces a number:
%
\begin{equation}
  g^{-1}(\omega, \sigma) = \inversemetric\,\omega_\mu\,\sigma_\nu \,.
  \label{eq:inverse-metric-bilinear}
\end{equation}
%
This is the inner product on the cotangent space, dual to the metric's inner
product on the tangent space.  We will use $\inversemetric$ repeatedly in later
chapters: in the geodesic Hamiltonian, it converts the covariant momentum
$p_\mu$ into the contravariant velocity $\dot{x}^\mu$.
\coderef{Math.Tensor.Operations}{inverseSym4}
\histnote{The index-manipulation calculus was developed by Gregorio
Ricci-Curbastro and Tullio Levi-Civita in 1900; Einstein adopted it for
general relativity and popularised the summation convention.}

\implnote{The function \texttt{inverseSym4} computes the inverse by cofactor
expansion, returning a \texttt{Sym4 'Contravariant}.  The variance flip from
\texttt{'Covariant} to \texttt{'Contravariant} is enforced by the type
signature --- passing a contravariant metric to this function produces a
compile-time error.}

\paragraph{A worked example: the two-sphere.}

Before tackling the full four-dimensional Lorentzian case, a two-dimensional
Riemannian example makes the mechanics of index operations concrete and
calculable by hand.  Consider the two-sphere $S^2$ of radius $R$, parametrised
by coordinates $(\theta, \varphi)$ with $\theta \in (0, \pi)$ and
$\varphi \in (0, 2\pi)$.  The metric is
%
\begin{equation}
  \d s^2 = R^2\,\d\theta^2 + R^2\sin^2\!\theta\;\d\varphi^2 \,,
  \label{eq:sphere-metric}
\end{equation}
%
so the metric components are $g_{\theta\theta} = R^2$,
$g_{\varphi\varphi} = R^2\sin^2\!\theta$, and $g_{\theta\varphi} = 0$.  As
with any coordinate chart on $S^2$, this one fails to cover the entire sphere
--- the determinant $g = R^4\sin^2\!\theta$ vanishes at the poles
$\theta = 0, \pi$, which is a coordinate singularity, not a geometric one.
The inverse metric has components $g^{\theta\theta} = 1/R^2$ and
$g^{\varphi\varphi} = 1/(R^2\sin^2\!\theta)$.

To see index raising in action, take a covector with components
$\omega_\theta = 0$, $\omega_\varphi = 1$.  Raising the index gives the
vector $\omega^\sharp$ with components
%
\begin{equation}
  \omega^\theta = g^{\theta\theta}\,\omega_\theta +
    g^{\theta\varphi}\,\omega_\varphi = 0 \,,\qquad
  \omega^\varphi = g^{\varphi\varphi}\,\omega_\varphi =
    \frac{1}{R^2\sin^2\!\theta} \,.
  \label{eq:sphere-raise-example}
\end{equation}
%
The raised component $\omega^\varphi$ depends on $\theta$: near the poles,
where $\sin\theta \to 0$, the same covector corresponds to a vector with
diverging contravariant component.  This reflects the coordinate singularity,
not any physical pathology --- the covector $\omega$ and the vector
$\omega^\sharp$ represent the same geometric object expressed in different
index positions.

\paragraph{From flat to curved.}

The Minkowski metric $\eta_{\mu\nu}$ from \cref{sec:sr-minkowski} is the
special case where the metric components are constants: the same diagonal
$(-1, +1, +1, +1)$ at every point.  In that setting, raising and lowering
merely flip the sign on the time component and leave spatial components
unchanged.  A general curved-spacetime metric $\metric$ can have all ten
components varying from point to point, encoding gravitational effects through
the geometry itself.
\xrefnote{The Schwarzschild metric, the first curved-spacetime metric we study
in detail, appears in \cref{ch:schwarzschild}.}

As a preview, the line element for a Schwarzschild black hole of mass $M$
reads
%
\begin{equation}
  \d s^2 = -\!\left(1 - \frac{2M}{r}\right)\d t^2
    + \left(1 - \frac{2M}{r}\right)^{\!-1}\!\d r^2
    + r^2\,\d\theta^2
    + r^2\sin^2\!\theta\;\d\varphi^2 \,,
  \label{eq:schwarzschild-preview}
\end{equation}
%
in coordinates $(t, r, \theta, \varphi)$ with geometric units $G = c = 1$.
The position-dependent coefficients encode the gravitational field: they
approach the Minkowski values at large $r$ and develop a coordinate singularity
at $r = 2M$ that coincides with the event horizon --- a genuine causal boundary
that requires horizon-penetrating coordinates to resolve.

The metric tensor connects the tangent-space framework to measurable physics.
Distances, proper times, angles, volumes, and the causal structure all follow
from $\metric$ and its inverse.  The next section introduces the derivative
operator that respects this structure: the covariant derivative, built from the
metric via the Christoffel symbols.
